\chapter{هم‌ترازی مغز، کدگذاری عصبی و استدلال}

\section{کدگذاری و رمزگشایی عصبی}

\subsection{\eng{Neural Encoding} و \eng{Neural Decoding}}
پس از پایان بحث شبکه‌های اسپایکی، مسیر بحث به سمت \eng{Neural Encoding} و \eng{Neural Decoding} رفت. هدف این بخش این است که روشن شود اطلاعات چگونه در مغز بازنمایی می‌شوند و چگونه می‌توان از روی فعالیت مغزی، رفتار، فرمان یا متن را بازسازی کرد.

در علوم اعصاب، \eng{Encoding} به این معناست که یک محرک خارجی وارد سیستم عصبی می‌شود و به نوعی فعالیت یا فرمان در مغز تبدیل می‌گردد. صدا، تصویر، لمس یا هر محرک بیرونی دیگر می‌تواند چنین مسیری را آغاز کند.

در مقابل، \eng{Decoding} یعنی از روی فعالیت‌های مغزی بتوانیم یک خروجی خارجی را بازسازی کنیم؛ برای مثال یک فرمان حرکتی، یک رفتار، یک متن یا حتی نیت فرد را حدس بزنیم. پس مسیر کلی را می‌توان چنین خلاصه کرد:

\[
\text{محرک خارجی} \longrightarrow \text{فعالیت مغزی} \longrightarrow \text{فرمان یا رفتار}
\]

در \eng{Encoding} از محرک به فعالیت مغزی می‌رسیم، و در \eng{Decoding} مسیر را برعکس طی می‌کنیم: از فعالیت مغزی تلاش می‌کنیم محرک، نیت یا فرمان را بازسازی کنیم.

\subsection{نظریه‌های اصلی \eng{Neural Coding}}
سه نظریهٔ اصلی دربارهٔ کدگذاری عصبی در این بحث به ترتیب تاریخی بررسی شدند:

\begin{enumerate}
  \item \eng{Grandmother Cell Theory}
  \item \eng{Population Coding}
  \item \eng{Sparse Coding}
\end{enumerate}

هیچ‌کدام از این نظریه‌ها در این بحث به‌عنوان پاسخ مطلق معرفی نشدند. هرکدام بخشی از شهود لازم برای فهم مغز و مدل‌های محاسباتی امروزی را فراهم می‌کنند.

\subsubsection{\eng{Grandmother Cell}}
نظریهٔ \eng{Grandmother Cell} حدود دههٔ ۱۹۶۰ مطرح شد. طبق این دیدگاه، برای هر مفهوم، شیء یا موجودیت خارجی، یک سلول یا یک گروه بسیار محدود و اختصاصی از نورون‌ها وجود دارد. مثال کلاسیک این است که برای مفهوم «مادربزرگ»، یک نورون یا گروهی خاص از نورون‌ها مسئول تشخیص همان مفهوم باشد.

نکتهٔ اصلی این نظریه انحصار است. منظور لزوماً این نیست که دقیقاً فقط یک نورون وجود دارد، بلکه ایده این است که مسئولیت شناخت یک مفهوم به‌شکل اختصاصی در اختیار یک سلول یا گروه کوچک قرار می‌گیرد.

این نظریه چند اشکال روشن دارد. نخست، اگر آن سلول یا گروه خاص آسیب ببیند، باید کل مفهوم از دست برود؛ درحالی‌که حافظه و شناخت معمولاً تا این اندازه به یک واحد منفرد وابسته نیستند. دوم، اگر برای هر چهره، صدا، شیء، مفهوم و حالت جهان به یک نورون یا گروه اختصاصی نیاز باشد، ظرفیت لازم بسیار بزرگ و از نظر زیستی و محاسباتی نامعقول می‌شود.

با این حال، ایدهٔ \eng{Grandmother Cell} کاملاً بی‌فایده نیست. در برخی سطوح مدل‌های عصبی، به‌ویژه در لایهٔ خروجی شبکه‌های طبقه‌بندی، شباهتی با این ایده دیده می‌شود. اگر یک شبکه سه کلاس داشته باشد و برای هر کلاس یک نود خروجی تعریف شود، هر نود تا حدی نقشی شبیه یک بازنمایی اختصاصی برای آن کلاس پیدا می‌کند. این تشبیه، نسخهٔ زیستی و مطلق \eng{Grandmother Cell} را اثبات نمی‌کند، اما نشان می‌دهد ایدهٔ «نمایندهٔ اختصاصی» در برخی معماری‌ها ظاهر می‌شود.

\subsubsection{\eng{Population Coding}}
\eng{Population Coding} در برابر دیدگاه اختصاصی \eng{Grandmother Cell} قرار می‌گیرد و حدوداً از دههٔ ۱۹۸۰ برجسته‌تر شد. در این نظریه، یک سلول یا گروه کوچک مسئول یک مفهوم نیستند؛ بلکه کل جمعیت نورون‌ها در تصمیم‌گیری نقش دارند.

این ایده را می‌توان به یک نظرسنجی جمعی تشبیه کرد. خروجی نهایی از مشارکت تعداد زیادی نورون ساخته می‌شود و هر نورون، بسته به فعالیت و وزن خود، سهمی در تصمیم دارد. در مدل‌های عصبی مصنوعی نیز شکل سادهٔ این ایده را در جمع وزن‌دار می‌بینیم:

\[
\sum_i W_i X_i
\]

در اینجا \(X_i\) می‌تواند فعالیت یا اسپایک نورون باشد و \(W_i\) وزن سیناپسی یا وزن یادگرفته‌شدهٔ آن. همهٔ نورون‌ها به یک اندازه تعیین‌کننده نیستند، اما تصمیم نهایی حاصل مشارکت جمعی آن‌هاست.

برای توضیح این نظریه، می‌توان به آزمایش‌هایی روی حرکت میمون‌ها اشاره کرد. در این نوع آزمایش‌ها، میمون برای حرکت یا پرتاب به جهت‌های مختلف مانند چپ، راست، بالا و پایین آموزش داده می‌شود و سپس فعالیت نورون‌های مغزی هنگام اجرای حرکت بررسی می‌گردد. مشاهدهٔ مهم این است که هنگام حرکت به یک جهت خاص، فقط یک نورون یا یک ناحیهٔ کاملاً اختصاصی فعال نمی‌شود؛ بلکه تعداد زیادی نورون با شدت‌های متفاوت درگیر می‌شوند. هنگام حرکت به چپ، نورون‌هایی که با جهت چپ سازگارترند فعالیت بیشتری دارند، اما نورون‌های دیگر هم می‌توانند با شدت کمتر فعال باشند.

مزیت نخست \eng{Population Coding} پایداری است. اگر یک نورون آسیب ببیند، کل مفهوم یا فرمان از بین نمی‌رود، چون اطلاعات در یک جمعیت بزرگ‌تر توزیع شده است. مزیت دوم دقت بالاتر است؛ چون تصمیم جمعی می‌تواند از تصمیم یک واحد منفرد پایدارتر باشد، هرچند این گزاره مطلق نیست. مزیت سوم ظرفیت محاسباتی بیشتر است: به‌جای اینکه برای هر زاویه، رنگ یا شیء یک نورون اختصاصی داشته باشیم، ترکیب فعالیت نورون‌ها می‌تواند حالت‌های بسیار متنوعی را بسازد.

ضعف اصلی این نظریه مصرف انرژی و تفسیرپذیری پایین‌تر است. اگر جمعیت بزرگی از نورون‌ها در بسیاری از تصمیم‌ها مشارکت کنند، هزینهٔ انرژی بالا می‌رود. از طرف دیگر، دیگر نمی‌توان گفت یک نورون مشخص مسئول یک خروجی مشخص است. به همین دلیل، \eng{Population Coding} از نظر تفسیرپذیری به یک \eng{black box} نزدیک می‌شود.

\subsubsection{\eng{Sparse Coding}}
\eng{Sparse Coding} را می‌توان حالتی میانی میان دو نظریهٔ قبلی دانست. این دیدگاه حدوداً از دههٔ ۱۹۹۰ پررنگ‌تر شد. در \eng{Sparse Coding}، نه همهٔ اطلاعات در اختیار یک سلول یا گروه بسیار کوچک است، و نه همهٔ نورون‌ها در همهٔ وظایف فعال می‌شوند. برای هر وظیفه، فقط زیرمجموعه‌ای نسبتاً کوچک و مرتبط از نورون‌ها فعال می‌شود.

در این حالت، ماتریس فعالیت نورونی تنک است. ممکن است برای یک کار خاص فقط ۲۰ درصد نورون‌ها، یا حتی ۵ درصد آن‌ها، فعال باشند و بقیه خاموش بمانند. این روش هم مصرف انرژی را کاهش می‌دهد و هم تفسیرپذیری را نسبت به \eng{Population Coding} بهتر می‌کند.

\eng{Sparse Coding} با مفهوم سلسله‌مراتب نیز ارتباط دارد. در پردازش تصویر، لایه‌های ابتدایی ممکن است لبه‌ها را تشخیص دهند، لایه‌های میانی رنگ، شکل یا ترکیب‌های پیچیده‌تر را بازنمایی کنند، و لایه‌های انتهایی به مفهوم یا شیء نهایی نزدیک شوند. در چنین معماری‌ای، گروه‌های مختلف نورونی در سطح‌های متفاوتی از انتزاع فعال می‌شوند.

مزیت \eng{Sparse Coding} این است که از یک طرف مانند \eng{Population Coding} به یک نورون منفرد وابسته نیست و از طرف دیگر مانند آن پرمصرف و کم‌تفسیر نیست. به همین دلیل، می‌توان آن را ترکیبی از پایداری نسبی و تفسیرپذیری نسبی دانست.

\subsection{ترکیب نظریه‌های کدگذاری در مدل‌های امروزی}
در مدل‌های محاسباتی امروز، بهتر است این سه نظریه را رقیب‌های مطلق ندانیم. در بسیاری از معماری‌ها می‌توان هر سه سطح را هم‌زمان دید:

\begin{align*}
\text{لایه‌های ابتدایی} &\longrightarrow \text{\eng{Population Coding}}\\
\text{لایه‌های میانی} &\longrightarrow \text{\eng{Sparse Coding}}\\
\text{لایه‌های انتهایی} &\longrightarrow \text{\eng{Grandmother-like Representation}}
\end{align*}

در لایه‌های ابتدایی، اطلاعات خام‌تر و کم‌انتزاع‌تر است. در پردازش صدا، این بخش می‌تواند به ویژگی‌هایی مانند واج، فرکانس یا الگوهای آوایی مربوط باشد. در بینایی نیز نواحی اولیه مانند \eng{V1} تا \eng{V4} ویژگی‌هایی مثل لبه، جهت، رنگ و فرم‌های ابتدایی را پردازش می‌کنند. این نمایش‌ها قدرتمندند، اما تفسیر هر بعد آن‌ها دشوار است؛ شبیه \eng{embedding}های مدل‌های یادگیری عمیق.

در لایه‌های میانی، ویژگی‌ها مشخص‌تر می‌شوند. ابتدا لبه‌ها و ویژگی‌های ساده‌تر دیده می‌شوند، سپس شکل‌ها و ترکیب‌ها و در ادامه ویژگی‌هایی که به اشیا یا مفاهیم نزدیک‌ترند. این سطح به \eng{Sparse Coding} نزدیک‌تر است، چون همهٔ واحدها برای همهٔ ورودی‌ها فعال نیستند.

در لایه‌های انتهایی، خروجی معمولاً به کلاس‌ها یا مفاهیم مشخص نگاشت می‌شود. اگر مدل چند کلاس داشته باشد، هر نود خروجی تا حدی مسئول یک کلاس است. این وضعیت شبیه یک بازنمایی \eng{Grandmother-like} است، با این تأکید که این شباهت به معنای پیاده‌سازی دقیق نظریهٔ زیستی \eng{Grandmother Cell} نیست.

\subsection{نسبت \eng{Symbolic AI} با نظریه‌های کدگذاری}
در این بحث این پرسش مطرح شد که اگر بخواهیم \eng{Symbolic AI} را با نظریه‌های کدگذاری عصبی مقایسه کنیم، به کدام نظریه نزدیک‌تر است. پاسخ اولیهٔ یکی از خوانندگان \eng{Grandmother Cell} بود، اما توضیح این بود که \eng{Symbolic AI} بیشتر به \eng{Sparse Coding} نزدیک است.

دلیل این مقایسه آن است که در \eng{Symbolic AI} معمولاً با سیستم‌های قانون‌محور، ویژگی‌های استخراج‌شده و ساختارهای نمادین سروکار داریم. تصمیم نهایی نه کاملاً توزیع‌شده و همه‌گیر مانند \eng{Population Coding} است، و نه کاملاً انحصاری مانند \eng{Grandmother Cell}. مجموعه‌ای از قواعد یا ویژگی‌های مشخص در تصمیم دخالت می‌کنند؛ بنابراین شباهت آن با \eng{Sparse Coding} بیشتر است.

\subsection{\eng{Brain-to-Text}}
بخش بعدی به مقاله‌ای از \eng{Meta} در سال ۲۰۲۵ دربارهٔ \eng{Brain-to-QWERTY} یا \eng{Brain-to-Text} اختصاص داشت. در چارچوب کلی \eng{Brain-to-Text} هدف این است که از فعالیت مغزی به متن برسیم، اما در \eng{Brain-to-QWERTY} دامنه محدودتر می‌شود: سیستم تلاش می‌کند بفهمد فرد هنگام تایپ چه کلیدهایی را روی صفحه‌کلید \eng{QWERTY} فشار داده یا قصد فشار دادن آن‌ها را داشته است.

پس هدف این مقاله خواندن تمام افکار فرد نیست. دامنهٔ مسئله به سناریوی تایپ محدود شده است، چون هرچه دامنه محدودتر باشد، \eng{decoding} فعالیت مغزی عملی‌تر می‌شود.

\subsubsection{مرحلهٔ آموزش}
در مرحلهٔ آموزش، فرد در دستگاه قرار می‌گیرد و فعالیت مغزی او هنگام خواندن و تایپ ثبت می‌شود. در این توضیح، روش‌هایی مانند \eng{EEG} و \eng{MEG} یا سیگنال‌های مشابه به‌عنوان مسیرهای ممکن ثبت فعالیت مغزی مطرح شدند.

روند کلی چنین است:

\begin{enumerate}
  \item متنی به فرد نشان داده می‌شود.
  \item فرد ابتدا متن را می‌خواند.
  \item سپس از او خواسته می‌شود متن را با صفحه‌کلید \eng{QWERTY} تایپ کند.
  \item هنگام تایپ، سیگنال‌های مغزی ثبت می‌شوند.
  \item سیستم می‌داند متن ورودی چه بوده و فرد قرار بوده چه چیزی تایپ کند.
  \item مدل یاد می‌گیرد فعالیت مغزی را به کلیدها یا متن تایپ‌شده نگاشت کند.
\end{enumerate}

در این مرحله، خروجی مدل با متن اصلی یا کلیدهای مورد انتظار مقایسه می‌شود و \eng{loss} محاسبه می‌گردد. سپس وزن‌های بخش‌های مختلف مدل بر اساس همین خطا به‌روزرسانی می‌شوند. این فرایند معمولاً برای همان فرد تکرار می‌شود تا مدل روی سیگنال‌های مغزی او کالیبره شود.

معماری کلی که در این بحث توصیف شد چنین بود:

\begin{align*}
\text{سیگنال مغزی هنگام تایپ} &\longrightarrow \text{\eng{CNN}}\\
&\longrightarrow \text{ویژگی‌های میانی}\\
&\longrightarrow \text{\eng{Transformer Decoder-only}}\\
&\longrightarrow \text{\eng{Language Model}}\\
&\longrightarrow \text{متن بازسازی‌شده}
\end{align*}

\coursefigure{0.95}{brain-to-qwerty-pipeline.png}{پایپ‌لاین \eng{Brain-to-QWERTY}: فرد جمله را می‌خواند و تایپ می‌کند، سیگنال‌های \eng{MEG/EEG} ثبت می‌شوند و مدل شامل ماژول کانولوشنی، ترنسفورمر و مدل زبان، متن هدف را بازسازی می‌کند.}

\subsubsection{نقش \eng{CNN}}
سیگنال مغزی خام ابتدا به یک \eng{Convolutional Neural Network} داده می‌شود. نقش \eng{CNN} استخراج ویژگی از سیگنال خام است. یعنی دادهٔ اولیهٔ مغزی به یک \eng{representation} یا مجموعه‌ای از ویژگی‌های میانی تبدیل می‌شود تا بخش‌های بعدی مدل بتوانند با آن کار کنند.

\subsubsection{نقش \eng{Transformer}}
ویژگی‌های استخراج‌شده به یک \eng{Transformer} داده می‌شوند. این \eng{Transformer} در این توضیح از نوع \eng{decoder-only} معرفی شد. وظیفهٔ آن مدل‌کردن توالی زبانی یا توالی کلیدهای تایپ‌شده از روی ویژگی‌های مغزی است؛ برای مثال تشخیص دهد فرد در حال تایپ چه حرف‌ها یا چه کلمه‌هایی بوده است.

\subsubsection{نقش \eng{Language Model}}
\eng{Language Model} در انتهای مسیر می‌تواند خطاهای خروجی را اصلاح کند. اگر بخش‌های قبلی مدل خروجی‌ای شبیه «صلام» تولید کنند، مدل زبانی با توجه به زمینه می‌تواند آن را به «سلام» نزدیک کند. بنابراین نقش آن شبیه یک تصحیح‌گر کانتکستی یا \eng{autocorrect} است.

\subsubsection{فاز آزمون}
در فاز آزمون، متن ورودی برای مدل معلوم نیست. فرد چیزی را تایپ می‌کند یا قصد تایپ آن را دارد و مدل فقط از روی فعالیت مغزی تلاش می‌کند بفهمد او چه نوشته یا قصد نوشتن چه چیزی داشته است. هدف نهایی این است که از سیگنال مغز بتوان به متن، فرمان یا نیت فرد رسید.

\subsection{کاربرد \eng{Brain-to-Text} و مقایسه با \eng{Neuralink}}
کاربرد مثبت چنین سیستم‌هایی کمک به افرادی است که ذهن فعال دارند اما توانایی حرکت یا ارتباط فیزیکی ندارند؛ برای مثال افرادی که نمی‌توانند تایپ کنند، دچار آسیب نخاعی‌اند یا محدودیت شدید حرکتی دارند. اگر فعالیت مغزی آن‌ها \eng{decode} شود، نیت یا دستور ذهنی می‌تواند به متن یا فرمان خارجی تبدیل شود و اجرای فرمان به سیستم‌های دیگر سپرده شود.

در همین زمینه، \eng{Neuralink} نیز نمونهٔ مهمی از ایمپلنت مغزی است که با جراحی در مغز قرار می‌گیرد و هدف آن تبدیل فعالیت مغزی، به‌ویژه نیت حرکتی، به فرمان خارجی است. برای مثال فرد به حرکت نشانگر ماوس فکر می‌کند و سیستم تلاش می‌کند همان نیت را به حرکت ماوس تبدیل کند.

تفاوت مهم این است که \eng{Brain-to-QWERTY} دامنهٔ محدودتری دارد و به تایپ روی صفحه‌کلید مربوط است، اما در \eng{Neuralink} هدف می‌تواند مستقیم‌تر باشد: تبدیل فکر یا نیت حرکتی به حرکت ماوس، کلیک یا کنترل یک ابزار.

\subsection{مقایسهٔ روش‌های ثبت فعالیت مغزی}
یکی از چالش‌های اصلی \eng{Brain-to-Text} محدودیت سخت‌افزار و روش ثبت فعالیت مغزی است. در این بحث سه دسته روش مقایسه شد: \eng{fMRI}، \eng{EEG} و روش‌های تهاجمی مانند ایمپلنت‌های مغزی.

\subsubsection{\eng{fMRI}}
\eng{fMRI} بر اساس تغییرات جریان خون کار می‌کند. وقتی بخشی از مغز فعال می‌شود، جریان خون در آن بخش تغییر می‌کند و \eng{fMRI} این تغییر را ثبت می‌کند.

مزیت \eng{fMRI} رزولوشن مکانی خوب آن است؛ یعنی می‌توان نسبتاً دقیق‌تر دید کدام ناحیهٔ مغز فعال شده است. ضعف آن رزولوشن زمانی پایین است، چون تغییرات جریان خون با تأخیر رخ می‌دهد. نکتهٔ مهم این است که \eng{fMRI} در مقیاس ثانیه معنا دارد، نه میلی‌ثانیه.

\subsubsection{\eng{EEG}}
\eng{EEG} فعالیت الکتریکی مغز را از روی پوست سر ثبت می‌کند. مزیت آن رزولوشن زمانی خوب است، چون تغییرات سریع الکتریکی را می‌تواند در حد میلی‌ثانیه ثبت کند. ضعف آن رزولوشن مکانی پایین‌تر است؛ سیگنال باید از بافت‌ها، جمجمه و پوست سر عبور کند و بنابراین تشخیص دقیق محل تولید سیگنال دشوارتر می‌شود.

\subsubsection{ایمپلنت‌های مغزی}
در روش‌های تهاجمی، حسگر یا تراشه با جراحی وارد مغز می‌شود. چون سنسور به نورون‌ها نزدیک‌تر است، هم رزولوشن زمانی خوب است و هم رزولوشن مکانی نسبت به روش‌های غیرتهاجمی بهتر می‌شود. اما مشکل اصلی، تهاجمی بودن این روش‌هاست: نیاز به جراحی، مسائل ایمنی، پایداری بلندمدت و ملاحظات اخلاقی.

بنابراین مقایسهٔ کلی چنین است:

\begin{itemize}
  \item \eng{fMRI}: غیرتهاجمی، رزولوشن مکانی خوب، رزولوشن زمانی پایین، مبتنی بر جریان خون.
  \item \eng{EEG}: غیرتهاجمی، رزولوشن زمانی خوب، رزولوشن مکانی پایین‌تر، مبتنی بر فعالیت الکتریکی از روی پوست سر.
  \item ایمپلنت‌های مغزی: تهاجمی، با رزولوشن زمانی و مکانی بهتر. نمونه‌ای مانند \eng{Neuralink} نیازمند جراحی است و چالش‌های پزشکی و اخلاقی دارد.
\end{itemize}

\subsection{چالش‌های \eng{Brain-to-Text}}
سه چالش اصلی که در این بحث برجسته شد عبارت‌اند از \eng{plasticity} مغز، تعمیم‌پذیری بین افراد و چندوجهی‌بودن فعالیت مغزی. این چالش‌ها باعث می‌شوند ساخت یک مدل عمومی و پایدار برای \eng{Brain-to-Text} بسیار دشوارتر از یک مسئلهٔ سادهٔ طبقه‌بندی باشد.

\subsubsection{\eng{Plasticity}}
مغز انسان ثابت نیست. هر فرد در طول زمان یاد می‌گیرد، فراموش می‌کند، خسته می‌شود و وضعیت ذهنی‌اش تغییر می‌کند. بنابراین مدلی که امروز روی سیگنال‌های مغزی یک فرد کالیبره شده، الزاماً در روزهای بعد دقیقاً همان ورودی‌ها را نمی‌بیند.

مسئله را می‌توان چنین توضیح داد: ما مدلی را روی مغز یک فرد آموزش می‌دهیم، اما خود آن مغز هم هم‌زمان در حال آموزش دیدن و تغییر کردن است. بنابراین مدل باید بتواند با \eng{plasticity} مغز کنار بیاید.

\subsubsection{تعمیم‌پذیری بین افراد}
نقشهٔ مغزی هر فرد با فرد دیگر متفاوت است و می‌توان آن را به اثر انگشت تشبیه کرد. مدلی که روی یک فرد آموزش دیده الزاماً روی فرد دیگر خوب کار نمی‌کند. در سامانه‌هایی مانند \eng{Neuralink} یا \eng{Brain-to-QWERTY} معمولاً مدل برای یک فرد خاص آموزش و آزمون می‌شود.

چالش اصلی این است که آیا می‌توان یک \eng{foundation model} برای مغز ساخت که روی افراد مختلف تعمیم پیدا کند یا نه. در مدل‌های زبانی تا حدی به چنین مدل‌هایی رسیده‌ایم، اما در مغز مسئله پیچیده‌تر است، چون باید فعالیت مغزی به زبان، حرکت، تصویر یا فرمان نگاشت شود.

\subsubsection{\eng{Multimodality}}
فعالیت مغزی ذاتاً چندوجهی است. حتی اگر آزمایش فقط روی متن تمرکز کند، مغز هم‌زمان از تصویر، صدا، حافظه، توجه، خستگی، احساسات و عوامل دیگر تأثیر می‌گیرد. به همین دلیل در آزمایش‌ها تلاش می‌شود دامنه محدود شود؛ برای مثال فقط متن نشان داده شود و فرد فقط تایپ کند.

اما هرچه سیستم به شرایط واقعی نزدیک‌تر شود، چندوجهی‌تر می‌شود و \eng{decoding} دشوارتر خواهد بود.

\subsection{مسیر آینده: سنسورهای پوشیدنی و سخت‌افزار \eng{Neuromorphic}}
مسیر ایده‌آل این حوزه حرکت به سمت سنسورهایی است که غیرتهاجمی، پوشیدنی، کم‌مصرف و دقیق باشند؛ یعنی هم رزولوشن زمانی خوب داشته باشند، هم رزولوشن مکانی مناسب، هم برای بدن قابل تحمل باشند و هم گرمای زیادی تولید نکنند. در این بحث، سنسورهای کوانتومی به‌عنوان نمونه‌ای از مسیرهای مورد توجه مطرح شدند.

مسئلهٔ انرژی در اینجا حیاتی است. اگر سخت‌افزاری نزدیک بدن یا داخل بدن قرار بگیرد، مصرف انرژی زیاد می‌تواند گرما تولید کند و برای بافت مغز خطرناک باشد. بنابراین سخت‌افزار باید هم دقیق و هم کم‌مصرف باشد.

یکی از ایده‌های مهم برای این هدف، سخت‌افزار \eng{neuromorphic} است. این سخت‌افزارها از ساختار و عملکرد مغز الهام می‌گیرند و معمولاً \eng{event-driven} هستند. در سیستم‌های \eng{time-based}، محاسبه بر اساس کلاک انجام می‌شود؛ چه رویدادی رخ داده باشد و چه رخ نداده باشد. اما در سیستم‌های \eng{event-driven}، فقط وقتی رویدادی اتفاق بیفتد پردازش انجام می‌شود. در شبکه‌های اسپایکی، این رویداد همان اسپایک است:

\[
\text{وقتی \eng{spike} رخ دهد} \longrightarrow \text{پردازش انجام می‌شود}
\]

چون مغز نیز تا حد زیادی رویدادمحور عمل می‌کند، سخت‌افزارهای \eng{neuromorphic} برای این حوزه گزینه‌ای طبیعی به نظر می‌رسند. این نکته به‌ویژه در \eng{Brain-to-Text} و ایمپلنت‌های مغزی مهم است.

\subsection{ترجیح مکانی، قانون \eng{Hebb} و مقداردهی اولیه}
در پایان این بخش، نکتهٔ مهم این است که در مدل‌های عصبی خروجی معمولاً از ترکیب وزن‌ها و فعالیت‌ها ساخته می‌شود:

\[
\sum_i W_i X_i
\]

اما این وزن‌ها فقط از یک عامل نمی‌آیند. سه عامل مهم در شکل‌گیری وزن نهایی مطرح شد: \eng{Hebbian Learning}، ترجیح مکانی و مقداردهی اولیه.

طبق قانون \eng{Hebb}، نورون‌هایی که با هم فعال می‌شوند، ارتباطشان تقویت می‌شود. این قانون در شکل‌گیری وزن‌های سیناپسی نقش دارد. علاوه بر آن، موقعیت مکانی یا ترجیح عملکردی نورون‌ها نیز مهم است. برای مثال، در آزمایش حرکت میمون، هنگام حرکت به سمت چپ، نورون‌هایی که از نظر مکانی یا عملکردی به جهت چپ نزدیک‌ترند بیشتر فعال می‌شوند.

عامل سوم مقداردهی اولیه است. افراد از ابتدا نیز تفاوت‌هایی در وزن‌ها یا تنظیمات اولیهٔ مغزی دارند و همین تفاوت‌ها در پاسخ نهایی نقش پیدا می‌کند. بنابراین وزن نهایی را می‌توان حاصل ترکیب این عوامل دانست:

\[
\text{وزن نهایی}
=
\text{\eng{Hebbian Learning}}
+
\text{ترجیح مکانی}
+
\text{مقداردهی اولیه}
\]

جمع‌بندی این بخش این است که سه نظریهٔ \eng{Grandmother Cell}، \eng{Population Coding} و \eng{Sparse Coding} سه نگاه متفاوت به بازنمایی عصبی ارائه می‌کنند. مدل‌های امروزی را می‌توان ترکیبی از این سه سطح دید: \eng{embedding}های اولیه به \eng{Population Coding} نزدیک‌اند، ویژگی‌ها و \eng{hidden state}های میانی به \eng{Sparse Coding}، و لایهٔ خروجی به بازنمایی \eng{Grandmother-like}. در ادامه، \eng{Brain-to-QWERTY} به‌عنوان یک نمونهٔ عملی از \eng{decoding} مغزی بررسی شد و چالش‌های آن، از تغییرپذیری مغز تا سخت‌افزار و انرژی، روشن شد.

\section{\eng{Brain Alignment} و \eng{Brain Score}}

\subsection{مسئلهٔ بازنمایی}
بحث \eng{Brain Alignment} پس از مباحث \eng{Brain-to-QWERTY} و \eng{Brain Decoding} مطرح شد. در بخش قبل دیدیم که می‌توان از روی سیگنال‌های مغزی، تا حدی قصد تایپ یا فرمان ذهنی فرد را بازسازی کرد. این کاربردها از یک طرف برای کمک به افراد دارای محدودیت حرکتی بسیار مهم‌اند، و از طرف دیگر نگرانی‌های جدی دربارهٔ ذهن‌خوانی، سوءاستفاده از داده‌های ذهنی و افشای ناخواستهٔ افکار ایجاد می‌کنند.

برای ورود به \eng{Brain Alignment}، مسئله را می‌توان از زاویهٔ بازنمایی توضیح داد. در یک نگاه فلسفی، می‌توان میان ماده، صورت یا \eng{Form}، و سطحی مجردتر تمایز گذاشت. آنچه در جهان مادی می‌بینیم، بازنمایی یا \eng{representation} چیزی عمیق‌تر است.

در همین چارچوب، خروجی یک مدل زبانی، سیگنال \eng{fMRI} هنگام فکر کردن، تصویر، متن، صدا یا رفتار انسان را می‌توان بازنمایی‌هایی متفاوت از یک معنا یا حقیقت واحد دانست. ممکن است چند بازنمایی متفاوت داشته باشیم، اما همهٔ آن‌ها به یک مفهوم مشترک اشاره کنند. مسئلهٔ اصلی این است که آیا بازنمایی‌های مدل و بازنمایی‌های مغز انسان به یکدیگر نزدیک‌اند یا نه.

\coursefigure{0.72}{platonic-representation-hypothesis.png}{فرضیهٔ بازنمایی افلاطونی: تصویر، متن و فعالیت مغزی را می‌توان به‌عنوان تصویرهایی متفاوت از یک واقعیت یا مفهوم زیرین مشترک در نظر گرفت.}

\subsection{آیا مدل واقعاً می‌فهمد؟}
نکتهٔ مهم این است که نمی‌توان به‌صورت صددرصدی و الگوریتمی ثابت کرد یک مدل «حقیقت» را فهمیده است یا نه. این مسئله تا حدی شبیه مسئلهٔ \eng{hallucination} است: هیچ الگوریتم قطعی‌ای وجود ندارد که تضمین کند مدل هیچ‌وقت دچار \eng{hallucination} نمی‌شود، چون درست یا غلط بودن پاسخ به زمینه وابسته است.

یک جمله ممکن است در زمینه‌ای تخیلی درست باشد، در زمینه‌ای پزشکی غلط باشد، یا از نگاه افراد مختلف تفسیرهای متفاوتی پیدا کند. بنابراین «فهمیدن» یا «نفهمیدن» مدل در معنای مطلق، مسئله‌ای کاملاً حل‌شده نیست.

با این حال، می‌توان بررسی کرد که مدل تا چه حد از پوستهٔ ظاهری داده فاصله گرفته و به معنا نزدیک‌تر شده است. مثال، سیب بود. اگر مدل فقط یاد گرفته باشد سیب معمولاً گرد، قرمز یا زرد است، ممکن است با دیدن یک سیب فاسد یا کپک‌زده دچار مشکل شود. اما انسان حتی اگر قبلاً سیب کپک‌زده ندیده باشد، مفهوم «فاسد شدن» را می‌فهمد و می‌تواند بگوید هنوز با یک سیب، اما یک سیب فاسد، روبه‌روست.

پس فهم انسانی صرفاً حفظ ظاهر نیست. مدل نیز اگر واقعاً به معنا نزدیک شده باشد، باید بتواند از ویژگی‌های سطحی عبور کند و ساختار انتزاعی‌تری از مفهوم را بازنمایی کند.

\subsection{ایدهٔ اصلی \eng{Brain Alignment}}
هدف \eng{Brain Alignment} این است که بررسی کنیم آیا بازنمایی‌های مدل با بازنمایی‌های مغز انسان هم‌راستا هستند یا نه. برای قضاوت دربارهٔ کیفیت یک مدل، فقط مقایسهٔ خروجی نهایی کافی نیست. باید ببینیم لایه‌های میانی مدل چه چیزی را \eng{encode} کرده‌اند، مغز انسان هنگام همان \eng{task} چه بازنمایی‌هایی دارد، و آیا این دو مجموعه بازنمایی قابل نگاشت به هم هستند یا نه.

اگر مدلی فقط خروجی درست بدهد، ممکن است شبیه طوطی الگوهایی را حفظ کرده باشد. اما اگر بازنمایی‌های میانی‌اش با بازنمایی‌های انسانی هم‌راستا باشند، احتمال بیشتری دارد که مدل به سطحی عمیق‌تر از معنا نزدیک شده باشد.

همچنین از روایت‌های تجربه‌های نزدیک به مرگ برای توضیح دشواری نگاشت میان سطح‌های مختلف بازنمایی استفاده کرد. در چنین روایت‌هایی، فرد گاهی تجربه‌ای غیرمادی را با مفاهیم مادی توضیح می‌دهد؛ مثلاً می‌گوید «درخت‌هایی بودند که حرف می‌زدند». یعنی تجربه را با واژگان جهان مادی بیان می‌کند، اما همان واژه‌ها ویژگی‌هایی پیدا می‌کنند که در مادهٔ عادی وجود ندارد. این مثال نشان می‌دهد نگاشت میان بازنمایی‌ها همیشه ساده و یک‌به‌یک نیست.

\subsection{طراحی آزمایش هم‌ترازی}
سؤال عملی این است: چگونه بفهمیم مدل از سطح ظاهر عبور کرده و به سطح معنا نزدیک شده است؟

مقایسهٔ مدل با انسان یک مسیر طبیعی است، اما اگر انسان را \eng{ground truth} بگیریم، باید مشخص کنیم چه چیزی را از انسان جمع‌آوری می‌کنیم. در بسیاری از آزمایش‌ها تا امروز، خروجی انسان و مدل مقایسه شده است: انسان \eng{task} را انجام می‌دهد، مدل هم همان \eng{task} را انجام می‌دهد، و خروجی‌ها مقایسه می‌شوند. اما برای \eng{Brain Alignment} فقط خروجی کافی نیست؛ بازنمایی‌های میانی نیز باید بررسی شوند.

یک مسیر دیگر استفاده از بازنمایی چندوجهی است. اگر حقیقت واحدی مانند مفهوم سیب از مسیرهای مختلفی مثل تصویر، متن، صدا و رفتار انسانی بررسی شود، و مدل بتواند در همهٔ این \eng{modality}ها همان مفهوم را تشخیص دهد، احتمالاً بازنمایی عمیق‌تری ساخته است.

مسیر سوم انتخاب \eng{task}هایی است که معنایی‌تر و دورتر از پوستهٔ مادی‌اند؛ مثل \eng{task}های فلسفی، استدلالی، علی و معلولی، فیزیکی، کشف قانون، استنتاج و قضاوت در وضعیت‌های جدید.

مثال نیوتون و افتادن سیب این پرسش را روشن می‌کند: اگر مدلی را زیر درخت بگذاریم و قبلاً قانون جاذبه را به آن نگفته باشیم، آیا می‌تواند از افتادن سیب تعجب کند و به وجود یک قانون برسد؟ در چنین \eng{task}هایی مطلوب است که مدل پاسخ را از قبل ندیده باشد، قانون به‌صورت صریح به آن داده نشده باشد، و مجبور شود رابطهٔ علی و معلولی بسازد.

هدف نهایی \eng{Brain Alignment} اثبات قطعی فهم مدل نیست. هدف عملی‌تر این است که بفهمیم مدل کجا را نمی‌فهمد تا بتوانیم آن را اصلاح کنیم. مسئلهٔ سخت «آیا مدل حقیقت را فهمیده است؟» به مسئلهٔ قابل‌بررسی‌تری کاهش داده می‌شود: دو جعبهٔ سیاه داریم، مغز انسان و مدل هوش مصنوعی؛ خروجی‌هایشان تا حدی مشابه است. حالا می‌خواهیم ببینیم لایه‌های میانی‌شان نیز با هم هم‌راستا هستند یا نه.

\subsubsection{تعریف لایه‌های میانی مغز}
در مدل عصبی مصنوعی، لایه‌ها و \eng{hidden state}ها مشخص‌اند. اما در مغز باید تصمیم بگیریم که «لایهٔ میانی» را چگونه تعریف کنیم. این پرسش به مسئلهٔ \eng{Brain Parcellation} می‌رسد: مغز را چگونه به بخش‌ها یا \eng{parcel}های مختلف تقسیم کنیم تا بتوان بازنمایی‌های آن را با مدل مقایسه کرد؟

\subsubsection{انتخاب \eng{Task}}
باید مشخص شود چه \eng{task}هایی برای سنجش هم‌ترازی مناسب‌اند. \eng{task}های زبانی، پیش‌بینی، استدلالی، معنایی و چندوجهی می‌توانند گزینه باشند. هرچه \eng{task} فقط به خروجی سطحی وابسته نباشد و مدل را مجبور به ساختن بازنمایی معنایی کند، برای بررسی \eng{alignment} مناسب‌تر است.

\subsubsection{تعریف \eng{Score}}
پس از جمع‌آوری داده از مغز و مدل، باید مشخص شود هم‌ترازی چگونه اندازه‌گیری می‌شود. معیار می‌تواند \eng{similarity}، \eng{correlation}، \eng{loss} یا ترکیبی از چند معیار باشد. مهم این است که یک معیار عملی داشته باشیم که نشان دهد بازنمایی مدل تا چه حد می‌تواند دادهٔ مغزی یا رفتاری انسان را پیش‌بینی کند.

\subsection{مقالهٔ \eng{Brain Score}}
\eng{Brain Score} به‌عنوان نمونه‌ای از این خط پژوهشی مطرح شد. تمرکز این مقاله بیشتر روی سؤال سوم است: چگونه برای \eng{alignment} یک \eng{benchmark} و \eng{score} تعریف کنیم؟

ایدهٔ کلی آزمایش چنین است: انسان داخل دستگاه \eng{fMRI} قرار می‌گیرد، یک محرک مثل متن به او داده می‌شود، همان متن به مدل نیز داده می‌شود، از مغز انسان سیگنال \eng{fMRI} گرفته می‌شود، و از مدل بازنمایی‌های میانی یا \eng{hidden vectors} استخراج می‌گردد. سپس یک مدل ساده یاد می‌گیرد بازنمایی مدل را به فعالیت مغزی نگاشت کند. در مرحلهٔ آزمون بررسی می‌شود این نگاشت تا چه حد می‌تواند فعالیت مغز را پیش‌بینی کند.

\subsubsection{اجزای آزمایش}
در این آزمایش دو جعبهٔ سیاه داریم: مغز انسان و مدل، مثلاً یک \eng{LLM}. ورودی یا \eng{stimulus} می‌تواند متن باشد.

برای مدل، متن وارد شبکه می‌شود و از لایه‌های میانی بردارهایی مانند \(v_1\)، \(v_2\) و \(v_3\) گرفته می‌شود. مدل اصلی \eng{freeze} می‌ماند، چون هدف ارزیابی آن است، نه آموزش دوبارهٔ آن.

برای انسان، همان متن نمایش داده می‌شود و از مغز او \eng{fMRI} گرفته می‌شود. ممکن است رفتارهایی مثل سرعت خواندن نیز ثبت شوند. سپس یک مدل کوچک و ساده آموزش داده می‌شود تا نگاشت میان بازنمایی‌های مدل و سیگنال‌های مغزی یا رفتاری را یاد بگیرد.

\subsubsection{مدل \eng{Mapping}}
مدل اصلی، مثلاً \eng{LLM}، دست‌کاری نمی‌شود. اما یک مدل کمکی آموزش داده می‌شود. وظیفهٔ این مدل کمکی این است که اگر بردار میانی مدل را دریافت کند، پیش‌بینی کند کدام الگوی \eng{fMRI} متناظر با آن است؛ یا اگر \eng{confidence} مدل در \eng{next token prediction} را دریافت کند، آن را به رفتار انسانی مثل سرعت خواندن نگاشت کند.

این مدل کمکی معمولاً می‌تواند خطی باشد. هدف ساختن یک مدل پیچیدهٔ جدید نیست؛ هدف این است که ببینیم آیا بازنمایی‌های موجود مدل با مغز قابل نگاشت هستند یا نه.

یکی از ایده‌های مهم، مقایسهٔ رفتار انسان با \eng{confidence} مدل است. اگر جمله برای انسان قابل پیش‌بینی و روان باشد، سریع‌تر خوانده می‌شود. اگر جمله غیرمنتظره باشد، انسان مکث می‌کند. برای مدل هم اگر \eng{next token} قابل پیش‌بینی باشد، \eng{confidence} بالاتر است؛ و اگر عبارت غیرمنتظره باشد، \eng{confidence} پایین‌تر می‌آید. بنابراین می‌توان \eng{confidence} مدل را با سرعت خواندن انسان مرتبط کرد.

نمونه این بود که «پردازش زبان طبیعی» عبارتی آشناست و ادامهٔ آن آسان‌تر پیش‌بینی می‌شود، اما «پردازش زبان فیزیک» عجیب‌تر است و می‌تواند باعث مکث شود.

\subsubsection{مرحلهٔ آموزش و آزمون}
در مرحلهٔ آموزش، متن‌هایی به انسان و مدل داده می‌شود. بازنمایی‌های مدل استخراج می‌گردد، \eng{fMRI} و رفتار انسان ثبت می‌شود، و مدل \eng{mapping} یاد می‌گیرد این‌ها را به هم وصل کند.

در مرحلهٔ آزمون، متن جدید به مدل داده می‌شود. مدل \eng{mapping} از روی بازنمایی مدل پیش‌بینی می‌کند که \eng{fMRI} یا رفتار انسانی چه خواهد بود. سپس این پیش‌بینی با دادهٔ واقعی انسان مقایسه می‌شود و از این مقایسه \eng{score} نهایی ساخته می‌شود.

\eng{Brain Score} می‌تواند از چند مؤلفه تشکیل شود؛ برای مثال هم‌ترازی با \eng{fMRI}، هم‌ترازی با خروجی یا پیش‌بینی، و هم‌ترازی با رفتارهایی مثل \eng{reading time}. در این توضیح، معیار اصلی برای مقایسه می‌تواند \eng{correlation} باشد.

\coursefigure{0.95}{brain-score-framework.png}{چارچوب \eng{Brain Score}: محرک زبانی مشترک به مدل و انسان داده می‌شود، سپس بازنمایی‌های داخلی مدل با پاسخ‌های مغزی و رفتاری انسان مقایسه می‌شوند.}

\subsubsection{اهمیت \eng{Next Word Prediction}}
در \eng{Brain Score} چند \eng{task} بررسی شده‌اند؛ از جمله \eng{question answering}، \eng{coherence}، \eng{entailment}، \eng{sentence similarity} و \eng{next word prediction}. نکتهٔ جالب این است که همبستگی مثبت و معنادار بیشتر در \eng{task} مربوط به \eng{next word prediction} دیده شده است.

این نتیجه با مباحث قبلی ترم سازگار است: مغز انسان دائماً در حال پیش‌بینی است. بنابراین \eng{task}هایی که حالت \eng{predictive} دارند، احتمالاً برای \eng{Brain Alignment} مناسب‌ترند. \eng{Task}هایی که صرفاً به خروجی نهایی مربوط‌اند، ممکن است بیشتر لایهٔ آخر را درگیر کنند؛ اما \eng{task}های پیش‌بینی‌کننده بازنمایی‌های میانی را بهتر آشکار می‌کنند.

\subsection{\eng{Brain Parcellation}}
\eng{Brain Parcellation} پاسخی به این پرسش است که لایه‌ها یا بخش‌های میانی مغز را چگونه تعریف کنیم. در این روش، مغز یا یک ناحیهٔ مورد علاقه از مغز به واحدهای کوچک‌تر و مجزا تقسیم می‌شود.

\subsubsection{\eng{Voxel}}
در \eng{fMRI}، مغز به واحدهای کوچک مکعبی تقسیم می‌شود که به هرکدام \eng{voxel} گفته می‌شود. \eng{Voxel}ها مرزهای طبیعی و زیستی مغز نیستند؛ آن‌ها واحدهایی‌اند که ما برای اندازه‌گیری تعریف می‌کنیم.

\subsubsection{\eng{Parcel}}
چند \eng{voxel} که در یک گروه قرار می‌گیرند، یک \eng{parcel} می‌سازند:

\[
\text{چند \eng{voxel}} \longrightarrow \text{یک \eng{parcel}}
\qquad
\text{چند \eng{parcel}} \longrightarrow \text{یک \eng{atlas} یا \eng{parcellation}}
\]

\subsubsection{\eng{Region of Interest}}
گاهی لازم نیست کل مغز بررسی شود. بسته به \eng{task}، فقط یک ناحیهٔ خاص انتخاب می‌شود. برای مثال اگر \eng{task} زبانی باشد، می‌توان فقط ناحیه‌های مرتبط با زبان را بررسی کرد. به چنین ناحیه‌ای \eng{Region of Interest} یا \eng{ROI} می‌گوییم.

\subsubsection{سیگنال \eng{BOLD}}
در \eng{fMRI} معمولاً از شاخصی به نام \eng{BOLD} استفاده می‌شود:

\[
\text{\eng{BOLD}} = \text{\eng{Blood Oxygen Level Dependent}}
\]

این شاخص تغییرات سطح اکسیژن خون را نشان می‌دهد. چون فعالیت عصبی با مصرف اکسیژن و جریان خون مرتبط است، تغییرات \eng{BOLD} به‌عنوان نشانه‌ای از فعالیت مغزی استفاده می‌شود.

برای هر \eng{voxel}، در طول زمان یک \eng{time series} داریم. اگر یک \eng{parcel} شامل چند \eng{voxel} باشد، می‌توان برای آن \eng{parcel} یک بردار میانگین یا بازنمایی کلی ساخت. برای مثال، اگر یک \eng{parcel} سه \eng{voxel} داشته باشد و در چهار زمان اندازه‌گیری انجام شود، دادهٔ اولیه یک ماتریس \(3 \times 4\) است که می‌توان آن را به بردار میانگین \eng{parcel} تبدیل کرد.

\subsection{روش‌های \eng{Parcellation}}
دو رویکرد کلی برای \eng{parcellation} مطرح شد: تقسیم‌بندی بر اساس آناتومی و تقسیم‌بندی بر اساس کارکرد.

\subsubsection{بر اساس آناتومی}
در این روش، مغز بر اساس ساختار آناتومیک تقسیم می‌شود. در بینایی، نمونه‌هایی مانند \eng{V1}، \eng{V2}، \eng{V4} و \eng{IT} مطرح‌اند. در اینجا تقسیم‌بندی بر اساس محل و ساختار فیزیکی مغز است.

\subsubsection{\eng{Functional Parcellation}}
در بحث \eng{Brain Alignment}، \eng{functional parcellation} معمولاً مهم‌تر است. در این روش، مهم نیست بخش‌ها از نظر فیزیکی دقیقاً کجا هستند؛ مهم این است که کدام بخش‌ها با هم کار می‌کنند. ممکن است دو نقطه از مغز از نظر مکانی دور باشند، اما چون هنگام یک کارکرد خاص با هم فعال می‌شوند، در یک شبکه قرار بگیرند.

این وضعیت را به کارمندانی تشبیه کرد که در طبقه‌های مختلف یک سازمان هستند، اما برای یک کار مشخص مدام با هم ارتباط دارند. از روی ارتباط کاری می‌فهمیم در یک تیم‌اند، نه صرفاً از روی محل فیزیکی‌شان.

\subsubsection{\eng{Localizer} و \eng{Atlas}}
برای انتخاب ناحیه‌های مغزی دو راه وجود دارد. راه اول استفاده از \eng{atlas}های موجود است؛ یعنی \eng{parcellation}هایی که قبلاً در مقاله‌های دیگر ساخته شده‌اند. پژوهشگر می‌گوید مبنای من فلان \eng{atlas} معروف است.

راه دوم ساختن \eng{localizer} مخصوص همان \eng{task} است. برای مثال اگر \eng{task} زبانی باشد، محرک زبانی به فرد داده می‌شود و دیده می‌شود کدام بخش‌های مغز فعال می‌شوند. سپس همان بخش‌ها بررسی می‌گردند. در این حالت پژوهشگر لزوماً از قبل نمی‌گوید ناحیهٔ زبان دقیقاً کجاست؛ بلکه بر اساس دادهٔ همان \eng{task} مشخص می‌کند کدام بخش‌ها فعال شده‌اند.

\subsection{شبکه‌های هفت‌گانهٔ \eng{Yeo}}
\eng{Yeo 7-Network Parcellation} یکی از تقسیم‌بندی‌های معروف کارکردی است. در این تقسیم‌بندی، مغز به هفت شبکهٔ کارکردی تقسیم می‌شود. این تقسیم‌بندی \eng{functional} است؛ یعنی بر اساس کارکرد و هم‌فعال‌شدن نواحی مغز ساخته شده، نه صرفاً آناتومی.

\coursefigure{0.82}{yeo-seven-network-parcellation.png}{شبکه‌های هفت‌گانهٔ \eng{Yeo}: افرازبندی کارکردی قشر مغز به شبکه‌های بینایی، حس‌حرکتی، پیش‌فرض، لیمبیک، توجه پشتی، توجه شکمی و فرونتوپاریتال.}

\subsubsection{\eng{Visual Network}}
\eng{Visual Network} مربوط به پردازش دیداری است. محل اصلی آن در بخش پس‌سری مغز است؛ یعنی تقریباً پشت سر، در راستای پردازش ورودی چشم‌ها.

\subsubsection{\eng{Somatomotor Network}}
\eng{Somatomotor Network} به حرکت و حس بدنی مربوط است؛ از جمله حرکت دست و پا، هماهنگی عضلات، انقباض عضلات و کنترل حرکتی بدن.

\subsubsection{\eng{Dorsal Attention Network}}
\eng{Dorsal Attention Network} به توجه درونی و ارادی مربوط است. وقتی خودمان تصمیم می‌گیریم روی چیزی تمرکز کنیم، مثلاً می‌خواهیم ساکت باشیم و متنی را بخوانیم، این شبکه درگیر می‌شود. این نوع توجه از درون هدایت می‌شود.

\subsubsection{\eng{Ventral Attention Network}}
\eng{Ventral Attention Network} به توجه بیرونی و ناگهانی مربوط است. وقتی کسی در می‌زند، صدایی ناگهانی می‌شنویم، یا چیزی از بیرون توجه ما را می‌گیرد، این شبکه فعال می‌شود. این توجه از تصمیم درونی نمی‌آید، بلکه با یک محرک بیرونی تحریک می‌شود.

برای توضیح \eng{dorsal} و \eng{ventral}، در این چارچوب می‌توان گفت اگر بدن یا مغز را در وضعیت جنینی تصور کنیم، \eng{dorsal} به سمت پشت یا بالا و \eng{ventral} به سمت شکم یا پایین اشاره دارد. در مغز نیز \eng{dorsal} معمولاً به بخش‌های بالاتر و پشتی‌تر و \eng{ventral} به بخش‌های پایین‌تر و جلویی‌تر اشاره می‌کند.

\subsubsection{\eng{Frontoparietal Network}}
\eng{Frontoparietal Network} در بخش‌هایی از ناحیهٔ پیشانی و آهیانه‌ای فعال است. کارکرد اصلی آن مدیریت و هماهنگی چند \eng{task}، تصمیم‌گیری، مرتب‌سازی و کنترل رفتار است. نام \eng{frontoparietal} به این معنا نیست که کل بخش پیشانی و کل بخش آهیانه‌ای درگیر است؛ بلکه فقط بخش‌هایی از این نواحی که در این کارکرد فعال می‌شوند، در شبکه قرار می‌گیرند.

\subsubsection{\eng{Default Mode Network}}
\eng{Default Mode Network} زمانی فعال‌تر می‌شود که \eng{task} مشخصی انجام نمی‌دهیم؛ مثل استراحت، تفکر آزاد، خیال‌پردازی، رؤیاپردازی یا ایده‌پردازی آزاد. گاهی وقتی از یک \eng{task} خسته می‌شویم و آن را رها می‌کنیم، ناگهان ایده‌های بهتری به ذهنمان می‌آید. این می‌تواند با فعال‌شدن \eng{Default Mode Network} مرتبط باشد، چون مغز خاموش نشده و وارد حالت پردازش آزادتر شده است.

\subsubsection{\eng{Limbic Network}}
\eng{Limbic Network} با احساسات، هیجان‌ها و پردازش عاطفی مرتبط است.

\subsection{فعالیت شبکه‌ها و \eng{Resting-State fMRI}}
نکتهٔ مهم این است که همهٔ شبکه‌های مغز نمی‌توانند هم‌زمان با حداکثر شدت فعال باشند. مغز درون جمجمه‌ای محدود قرار دارد و جریان خون نیز محدود است. فعالیت مغزی شبیه پینگ‌پنگ میان شبکه‌هاست: بعضی شبکه‌ها فعال‌تر می‌شوند، بعضی پایین می‌آیند، و بسته به \eng{task} الگوی فعالیت تغییر می‌کند. برای مثال نمی‌توان هم‌زمان هم درگیر تفکر آزاد بود و هم با تمرکز کامل یک \eng{task} مشخص را انجام داد.

برای توضیح محدودیت فیزیکی مغز، می‌توان به جراحی مغز اشاره کرد. در برخی جراحی‌ها، بعد از باز کردن جمجمه، ممکن است نتوان بلافاصله استخوان را مثل قبل سر جای خود گذاشت، چون مغز دچار تورم می‌شود. این مثال نشان می‌دهد محدودیت فیزیکی جمجمه و جریان خون در فعالیت مغز نقش واقعی دارند.

در \eng{parcellation} معروف \eng{Yeo} از \eng{resting-state fMRI} استفاده شده است. در \eng{resting-state fMRI} فرد \eng{task} خاصی انجام نمی‌دهد، اما این به معنای خاموش بودن مغز نیست. مغز حتی هنگام استراحت یا خواب فعال است، انرژی زیادی مصرف می‌کند، همچنان پیش‌بینی انجام می‌دهد و شبکه‌های آن کاملاً خاموش نمی‌شوند.

از نظر انرژی نیز همین نکته مهم است: مغز انسان با اینکه فقط بخش کوچکی از وزن بدن را تشکیل می‌دهد، حدود ۲۰ درصد انرژی بدن را مصرف می‌کند. انجام کار شناختی سخت می‌تواند الگوی فعالیت شبکه‌ها را عوض کند، اما افزایش مصرف انرژی نسبت به فعالیت پایه معمولاً کوچک است و در حد چند درصد گزارش می‌شود. بنابراین \eng{resting state} حالت خاموشی نیست؛ بیشتر یک خط پایهٔ فعال و پرهزینه است که \eng{task}ها روی آن تغییرات نسبی ایجاد می‌کنند.

\subsection{همبستگی در \eng{Parcellation}}
برای \eng{parcellation}، \eng{time series} سیگنال \eng{BOLD} وکسل‌ها بررسی می‌شود. اگر دو \eng{voxel} یا دو ناحیهٔ مغزی در طول زمان با هم بالا و پایین شوند، یعنی همبستگی مثبت بالایی داشته باشند، احتمالاً در یک شبکهٔ کارکردی قرار دارند.

روند ساده چنین است:

\begin{enumerate}
  \item برای هر \eng{voxel} یک \eng{time series} گرفته می‌شود.
  \item همبستگی میان \eng{time series}ها محاسبه می‌شود.
  \item \eng{voxel}هایی که همبستگی بالایی دارند \eng{cluster} می‌شوند.
  \item هر \eng{cluster} به‌عنوان یک \eng{parcel} یا \eng{network} در نظر گرفته می‌شود.
\end{enumerate}

در اینجا \eng{clustering} همان \eng{parcellation} است.

\subsection{\eng{Task-Positive} و \eng{Task-Negative}}
در پایان بحث \eng{parcellation}، شبکه‌ها را می‌توان با دو مفهوم توضیح داد. \eng{Task-positive networks} هنگام انجام یک \eng{task} فعال‌تر می‌شوند. شبکه‌های \eng{Dorsal Attention}، \eng{Ventral Attention} و \eng{Frontoparietal} معمولاً در این دسته‌اند.

در مقابل، \eng{task-negative networks} شبکه‌هایی‌اند که هنگام انجام \eng{task} کاهش فعالیت دارند و در حالت عدم انجام \eng{task} فعال‌ترند. این دسته به‌ویژه با \eng{Default Mode Network} مرتبط است.

\coursefigure{0.82}{task-positive-and-task-negative-networks.png}{نمونه‌ای از شبکه‌های وظیفه‌مثبت و وظیفه‌منفی و نیز شبکه‌های توجه پشتی و شکمی؛ شکل رابطهٔ رقابتی و همپوشانی محدود شبکه‌های کارکردی را نشان می‌دهد.}

\subsection{بازگشت به \eng{Brain Alignment}}
پس از \eng{parcellation} می‌توان بازنمایی‌های مغزی را با بازنمایی‌های مدل مقایسه کرد. روند کلی چنین است:

\begin{enumerate}
  \item مغز به \eng{parcel}های کارکردی تقسیم می‌شود.
  \item برای هر \eng{parcel} سیگنال \eng{BOLD} یا بازنمایی زمانی استخراج می‌شود.
  \item از مدل، \eng{hidden state}ها یا بردارهای لایه‌های میانی گرفته می‌شود.
  \item میان بازنمایی‌های مدل و مغز \eng{correlation} یا \eng{similarity} محاسبه می‌شود.
  \item هرچه این ارتباط بیشتر باشد، می‌توان گفت هم‌ترازی بیشتر است.
\end{enumerate}

با این حال، این همچنان اثبات فهم حقیقی نیست؛ فقط راهی عملی برای نزدیک‌شدن به مسئله است.

\subsection{خط پژوهشی آینده: استنتاج و مغالطه}
در پایان، در این چارچوب می‌توان گفت اگر بخواهیم واقعاً بررسی کنیم مدل به معنا و فهم انسانی نزدیک شده یا نه، باید از \eng{task}های سخت‌تر استفاده کنیم: طراحی \eng{task}های استنتاجی، ساختن شرایطی که مدل در آن به اشتباه بیفتد، بررسی مغالطه‌ها، و سنجش \eng{reasoning} انسانی و خطاهای انسانی در مدل.

مغالطه پدیده‌ای انسانی است. اگر مدلی بتواند مغالطه را تشخیص دهد یا حتی ساختار مغالطه را بازتولید کند، یعنی تا حدی ساختار استدلال را فهمیده است. این مسیر یعنی چالش‌کردن مدل از نظر استنتاجی، ادامهٔ طبیعی بحث هم‌ترازی و استدلال است.

\section{استدلال، قیاس و هم‌ترازی خطای انسان و مدل}

\subsection{چرا پاسخ درست کافی نیست؟}
بخش قبل نشان داد که در پس داده‌ها می‌توان یک حقیقت یا مفهوم پنهان را در نظر گرفت؛ چیزی که در این بحث با \(Z\) نشان داده شد. خود \(Z\) را مستقیماً در اختیار نداریم، اما بازنمایی‌هایی از آن داریم: بازنمایی در مغز انسان، بازنمایی در لایه‌های میانی مدل، و خروجی نهایی مدل.

در \eng{Brain Alignment} هدف این است که بفهمیم آیا مدل واقعاً به آن مفهوم پنهان نزدیک شده است یا فقط از نظر آماری خروجی‌هایی شبیه داده‌های آموزشی تولید می‌کند. چون \(Z\) را مستقیم نداریم، مغز انسان به‌عنوان مرجع عملی استفاده می‌شود. \eng{Brain Score} نیز در همین چارچوب می‌پرسد لایه‌های میانی مدل تا چه حد با لایه‌های میانی مغز انسان هم‌بستگی دارند.

اما نکتهٔ مهم این بخش این بود که هم‌ترازی نباید فقط بر اساس پاسخ درست سنجیده شود. اگر انسان‌ها در یک موقعیت خاص خطا می‌کنند، مدلی که واقعاً شبیه انسان می‌فهمد باید در الگوی خطا نیز تا حدی شبیه انسان باشد. بنابراین سؤال فقط این نیست که مدل چند درصد پاسخ درست داده است؛ سؤال‌های مهم‌تر این‌اند:

\begin{itemize}
  \item آیا مدل همان جاهایی اشتباه می‌کند که انسان‌ها اشتباه می‌کنند؟
  \item آیا خطاهای شایع انسانی در مدل نیز شایع‌اند؟
  \item آیا ترجیحات، زمان پاسخ، و توزیع خطاهای مدل به انسان شبیه است؟
\end{itemize}

پس هم‌ترازی می‌تواند چند سطح داشته باشد: هم‌بستگی لایه‌های میانی مدل با مغز، هم‌بستگی خروجی مدل با خروجی انسان، شباهت در زمان پاسخ، شباهت در ترجیحات، و شباهت در خطاها و بایاس‌های انسانی.

\subsection{یادگیری از خطای انسان}
در نگاه اول شاید عجیب باشد که بخواهیم مدل خطاهای انسان را نیز یاد بگیرد. در برخی کاربردها، مثل جراحی رباتیک، هدف دقیقاً این است که مدل خطای انسان را پوشش دهد و بهتر از انسان عمل کند. اما موضوع این بخش ساختن چنین سیستم‌هایی نبود؛ موضوع این بود که آیا مدل واقعاً مفهوم پشت داده را فهمیده است یا نه.

اگر مغز انسان نزدیک‌ترین مرجع ما به فهم مفهوم باشد، خطاهای انسانی نیز بخشی از همان مرجع می‌شوند. از این زاویه، سوگیری شناختی الزاماً باگ نیست. بسیاری از این خطاها در طول تکامل برای بقا مفید بوده‌اند.

برای مثال، انسانی که در موقعیت خطرناک بیش از حد صبر کند تا همه‌چیز را منطقی اثبات کند، ممکن است زنده نماند. در مقابل، تصمیم‌گیری سریع بر اساس پیش‌فرض‌های ذهنی می‌تواند شانس بقا را افزایش دهد. گاهی خطا داشتن یعنی ساده‌سازی و تصمیم‌گیری سریع‌تر.

نمونهٔ دیگر مصرف انرژی مغز است. مغز نمی‌خواهد همهٔ اطلاعات محیط را با جزئیات کامل پردازش کند؛ بنابراین فیلتر می‌گذارد و فقط بخشی از اطلاعات را برجسته می‌کند. وقتی فرد به یک محیط آموزشی نگاه می‌کند، به همهٔ صندلی‌ها و جزئیات محیط توجه نمی‌کند، بلکه بیشتر روی چهره‌ها، واکنش‌ها و صداها تمرکز می‌کند. این فیلترکردن از یک نظر خطاست، اما برای عملکرد طبیعی شناخت لازم است.

حتی طیف بینایی انسان نیز همین منطق را نشان می‌دهد. انسان فقط نور مرئی را می‌بیند، نه کل طیف الکترومغناطیسی. اگر چشم انسان مادون قرمز را نیز می‌دید، شاید دائماً جریان گرما، تغییرات دمای بدن افراد و حرکت هوای گرم و سرد را می‌دید و توجهش مختل می‌شد. حذف برخی اطلاعات ضعف مطلق نیست؛ نوعی بهینه‌سازی است.

\subsection{\eng{Pareidolia} و تشخیص الگو}
\eng{Pareidolia} یعنی دیدن چهره یا الگوی آشنا در چیزهایی که واقعاً چنین الگوی معناداری ندارند؛ مثل ابرها، نقش مبل، گلدوزی یا اشیای بی‌جان. برای مثال، کودکان خردسال گاهی روی نقش مبل یا اشیای دیگر چهره می‌بینند و آن را به انسان نسبت می‌دهند.

این رفتار الزاماً آموخته‌شده نیست، بلکه بخشی از سیستم ادراکی انسان است. از نظر منطقی، دیدن چهره در نقش بی‌جان خطاست؛ اما از نظر تکاملی می‌تواند مفید باشد، چون انسان باید چهره‌ها، نگاه‌ها و واکنش‌های دیگران را سریع تشخیص دهد. اگر کسی وارد یک جمع شود و اصلاً به چهرهٔ افراد واکنش نشان ندهد، برای ما ناخوشایند و غیرطبیعی به نظر می‌رسد.

پس \eng{Pareidolia} نمونه‌ای از خطای ادراکی است که در سطح بقا و تعامل اجتماعی می‌تواند معنا داشته باشد.

\subsection{\eng{System 1}، \eng{System 2} و تفکر انتقادی}
در بحث استدلال دوباره به تفاوت \eng{System 1} و \eng{System 2} اشاره شد. \eng{System 1} سریع، کم‌هزینه و مبتنی بر \eng{prior knowledge} یا دانش قبلی است. این سیستم بر اساس تجربه‌های پیشین و الگوهای آماری سریع تصمیم می‌گیرد.

\eng{System 2} کندتر، پرهزینه‌تر و تحلیلی‌تر است. وقتی فعال می‌شود، تلاش می‌کند مسئله را واقعاً بررسی و اثبات کند. تفکر انتقادی بیشتر به \eng{System 2} مربوط است؛ یعنی فرد در موقعیت‌های لازم می‌تواند از پاسخ سریع و پیش‌فرض ذهنی عبور کند و وارد تحلیل عمیق‌تر شود.

\subsection{انواع \eng{Reasoning}}
\eng{Reasoning} یا استدلال انواع مختلفی دارد. در این بحث سه نوع اصلی نام برده شد:

\begin{itemize}
  \item \eng{Deductive Reasoning}
  \item \eng{Inductive Reasoning}
  \item \eng{Abductive Reasoning}
\end{itemize}

تمرکز این بخش روی شکل خاصی از \eng{Deductive Reasoning} بود: \eng{Syllogistic Reasoning}. در فارسی می‌توان آن را قیاس یا «صغرا و کبرا چیدن» نامید.

\subsection{\eng{Syllogism}}
\eng{Syllogism} ساختاری صوری برای استدلال قیاسی است. در ساده‌ترین حالت، دو مقدمه و یک نتیجه داریم:

\begin{enumerate}
  \item \eng{Premise 1}
  \item \eng{Premise 2}
  \item \eng{Conclusion}
\end{enumerate}

\subsubsection{ساختار قیاس}
در این ساختار یک \eng{Middle Term} یا ترم میانی وجود دارد که در هر دو مقدمه مشترک است و اجازه می‌دهد میان دو گزاره ارتباط برقرار کنیم. اگر ترم مشترک وجود نداشته باشد، اصولاً قیاسی شکل نمی‌گیرد، چون چیزی وجود ندارد که دو مقدمه را به هم وصل کند.

\subsubsection{\eng{Mode A}}
\eng{Mode A} رابطهٔ کلی مثبت است:

\[
\text{\eng{All A are B}}
\]

یعنی همهٔ \(A\)ها \(B\) هستند.

\subsubsection{\eng{Mode I}}
\eng{Mode I} رابطهٔ جزئی مثبت است:

\[
\text{\eng{Some A are B}}
\]

یعنی بعضی از \(A\)ها \(B\) هستند.

\subsubsection{\eng{Mode E}}
\eng{Mode E} رابطهٔ کلی منفی است:

\[
\text{\eng{No A are B}}
\]

یعنی هیچ \(A\)ای \(B\) نیست.

\subsubsection{\eng{Mode O}}
\eng{Mode O} رابطهٔ جزئی منفی است:

\[
\text{\eng{Some A are not B}}
\]

یعنی بعضی از \(A\)ها \(B\) نیستند.

\subsubsection{\eng{Figure}}
علاوه بر \eng{mode}ها، در \eng{Syllogism} چند \eng{figure} مختلف نیز داریم. \eng{Figure} یعنی جایگاه ترم‌ها، به‌ویژه ترم میانی، در دو مقدمه تغییر کند. با چهار \eng{mode} و چهار \eng{figure}، در مجموع \(4 \times 4 \times 4 = 64\) فرم ممکن برای \eng{Syllogism} ساخته می‌شود. بعضی از این فرم‌ها نتیجهٔ معتبر دارند و بعضی ندارند.

در صورت‌بندی رایج مسئله‌های قیاسی، از این ۶۴ فرم، ۲۷ فرم نتیجهٔ معتبر دارند و ۳۷ فرم به \eng{No Valid Conclusion} یا \eng{NVC} می‌رسند. بنابراین ارزیابی مدل فقط این نیست که در چند فرم به پاسخ درست برسد؛ باید دید آیا می‌تواند فرم‌هایی را که اساساً نتیجهٔ معتبر ندارند نیز از فرم‌های معتبر جدا کند یا نه.

چهار شکل اصلی را می‌توان به‌صورت زیر خلاصه کرد:
\[
\begin{array}{c@{\qquad}c@{\qquad}c@{\qquad}c}
\text{\eng{Figure 1}} & \text{\eng{Figure 2}} & \text{\eng{Figure 3}} & \text{\eng{Figure 4}}\\
A-B & B-A & A-B & B-A\\
B-C & C-B & C-B & B-C
\end{array}
\]

\subsubsection{\eng{Valid Conclusion}}
یک نتیجه زمانی \eng{valid} است که در همهٔ حالت‌های ممکن درست باشد و هیچ مثال نقضی نداشته باشد. فرم معروف \eng{Barbara} یا \((AAA)-1\) از ساده‌ترین فرم‌های معتبر است:

\begin{align*}
\text{\eng{All A are B}}\\
\text{\eng{All B are C}}\\
\therefore \quad \text{\eng{All A are C}}
\end{align*}

این نتیجه همیشه معتبر است. اما بسیاری از فرم‌ها نتیجهٔ معتبر ندارند. ممکن است در بعضی حالت‌ها یک نتیجه درست از آب دربیاید، اما اگر همیشه درست نباشد، \eng{valid} محسوب نمی‌شود. چنین مواردی می‌توانند منشأ مغالطه باشند.

\subsection{استفاده از قیاس برای \eng{Brain Alignment}}
پرسش اصلی این است که ساختارهای قیاسی چه ربطی به \eng{Brain Alignment} و مدل‌های زبانی دارند. می‌توان ۶۴ فرم \eng{Syllogism} را به انسان‌ها داد و دید در کدام فرم‌ها درست پاسخ می‌دهند و در کدام فرم‌ها خطا می‌کنند. سپس همان فرم‌ها را به مدل زبانی داد و خروجی مدل را با انسان مقایسه کرد.

در اینجا دقت کلی کافی نیست. مهم‌تر این است که توزیع خطاهای مدل شبیه توزیع خطاهای انسان باشد. اگر انسان‌ها در یک فرم خاص زیاد اشتباه می‌کنند، ولی مدل همیشه درست جواب می‌دهد، این الزاماً نشانهٔ هم‌ترازی نیست. ممکن است مدل فقط ساختارهای صوری را از دادهٔ آموزشی یاد گرفته باشد، بدون اینکه شبیه انسان استدلال کند.

\subsection{نتایج انسانی در \eng{Syllogistic Reasoning}}
در یک مطالعهٔ مرتبط، همین ۶۴ فرم \eng{Syllogism} به انسان‌ها داده شده بود تا مشخص شود انسان‌ها در فرم‌های دارای نتیجهٔ معتبر و فرم‌های \eng{NVC} چه الگوی پاسخی دارند.

\coursefigure{0.95}{analogy-accuracy-by-model-scale.png}{دقت مدل‌ها در فرم‌های مختلف قیاس با مقیاس‌های گوناگون؛ نمودار نشان می‌دهد افزایش اندازهٔ مدل می‌تواند دقت را از خط پایهٔ انسانی بالاتر ببرد.}

برای هر فرم، توزیع پاسخ انسان‌ها بررسی شده بود. حتی در فرم‌های ساده‌ای مثل \eng{Barbara} هم همهٔ افراد پاسخ درست نمی‌دهند: تعداد زیادی درست جواب می‌دهند، اما همچنان بخشی از افراد پاسخ اشتباه تولید می‌کنند. این نشان می‌دهد که حتی در قیاس‌های صوری ساده نیز انسان‌ها همیشه کاملاً منطقی عمل نمی‌کنند.

در این بحث این پرسش مطرح شد که چرا با وجود شناخته‌شدن این فرم‌ها در طول قرن‌ها، انسان‌ها همچنان در بسیاری از آن‌ها اشتباه می‌کنند. چند پاسخ ممکن مطرح شد: انسان در زندگی روزمره به همهٔ این فرم‌ها نیاز نداشته، برخی فرم‌ها کمتر در تجربهٔ واقعی استفاده شده‌اند، بعضی فرم‌های غلط یا مغالطه‌آمیز ممکن است در تعاملات اجتماعی کاربرد داشته باشند، و مغز انسان برای حل همهٔ فرم‌های منطقی بهینه نشده، بلکه برای تصمیم‌گیری سریع در جهان واقعی بهینه شده است.

بنابراین ممکن است انسان بتواند با تمرین همهٔ این فرم‌ها را یاد بگیرد، اما مغز به‌صورت طبیعی خودش را درگیر تمام حالت‌های ممکن نکرده است.

در یک فرم سخت‌تر، چند پاسخ مختلف ممکن است پیشنهاد شود، اما پاسخ معتبر ابتدا به‌راحتی پیدا نشد. نمودار ون نشان می‌دهد که برای \eng{valid} بودن، نتیجه باید در همهٔ حالت‌ها درست باشد. اگر فقط در بعضی حالت‌ها درست باشد، نتیجه معتبر نیست.

در \eng{Syllogism}های صوری، معمولاً بحث مجموعهٔ تهی وارد نمی‌شود، چون هدف بررسی قیاس رسمی است، نه همهٔ پیچیدگی‌های زبان طبیعی یا هستی‌شناسی. اما وقتی این ساختارها به زبان طبیعی برده می‌شوند، مسائلی مانند تهی بودن یک مجموعه می‌توانند مهم شوند.

\subsection{از فرم صوری به زبان طبیعی}
در حالت صوری، جمله‌ها به شکل‌هایی مانند \(\text{\eng{All A are B}}\) و \(\text{\eng{Some B are C}}\) نوشته می‌شوند. اما برای مقایسهٔ واقعی با انسان و مدل زبانی، بهتر است این‌ها به \eng{Natural Language} تبدیل شوند؛ یعنی به‌جای \(A\)، \(B\) و \(C\)، جمله‌های طبیعی و معنادار داشته باشیم.

در این حالت کار سخت‌تر می‌شود، چون مدل باید ترم میانی را تشخیص دهد، رابطهٔ منطقی میان دو مقدمه را پیدا کند، مغالطه‌های زبانی را تشخیص دهد، و نتیجهٔ معتبر یا نبود نتیجهٔ معتبر را تولید کند.

\subsection{دیتاست \eng{Avicenna}}
دیتاستی به نام \eng{Avicenna} حدود ۶۰۰۰ نمونهٔ \eng{Syllogism} در زبان طبیعی دارد.

در هر نمونه معمولاً این اجزا وجود دارد:

\begin{itemize}
  \item \eng{Premise 1}
  \item \eng{Premise 2}
  \item \eng{Syllogistic Relation}
  \item \eng{Conclusion} یا \eng{No Valid Conclusion}
\end{itemize}

\eng{Syllogistic Relation} مشخص می‌کند آیا واقعاً بین دو مقدمه ترم میانی معناداری وجود دارد یا نه. گاهی دو جمله از نظر واژگانی کلمات مشترک دارند، اما این اشتراک برای استدلال کافی نیست. مثلاً ممکن است در یک جمله «شیر جنگل» و در جملهٔ دیگر «شیر آب» آمده باشد. واژه مشترک است، اما ترم میانی واقعی وجود ندارد؛ پس استدلال معتبر شکل نمی‌گیرد.

\subsection{تفاوت \eng{Syllogism} و \eng{Entailment}}
در \eng{Entailment} ممکن است فقط یک گزاره داشته باشیم و با کمک دانش عمومی نتیجه‌ای بگیریم. برای مثال، اگر جمله‌ای دربارهٔ سیگار و آرسنیک باشد، ممکن است با دانش قبلی بدانیم آرسنیک ماده‌ای خاص است و از آن نتیجه‌ای بگیریم.

اما در \eng{Syllogism}، دو مقدمه باید به‌واسطهٔ یک ترم میانی به هم وصل شوند و نتیجه از ترکیب آن دو مقدمه به دست بیاید. پس هر نتیجه‌گیری زبانی الزاماً \eng{Syllogism} نیست.

\subsection{مغالطه‌ها و استدلال‌های روزمره}
برای نشان دادن نقش \eng{prior knowledge} و احساسات در استدلال، دو مثال روزمره مطرح شد.

مثال اول:

\begin{enumerate}
  \item دولت‌های فاسد تورم ایجاد می‌کنند.
  \item کشور ایران تورم دارد.
  \item پس دولت ایران فاسد است.
\end{enumerate}

این نتیجه از نظر منطقی معتبر نیست، چون تورم می‌تواند علت‌های دیگری هم داشته باشد. این ساختار شبیه مغالطهٔ تأیید تالی است. اما اگر فرد از قبل نسبت به دولت نارضایتی داشته باشد، نتیجه برای او پذیرفتنی به نظر می‌رسد. در اینجا \eng{prior knowledge} و احساسات قبلی باعث می‌شوند \eng{System 1} فعال شود و فرد وارد بررسی دقیق \eng{System 2} نشود.

مثال دوم:

\begin{enumerate}
  \item هر کشوری که تولید ناخالص داخلی‌اش رشد کند، اقتصادش بزرگ‌تر شده است.
  \item طبق آمار جهانی، تولید ناخالص داخلی ایران رشد کرده است.
  \item پس اقتصاد ایران بزرگ‌تر شده است.
\end{enumerate}

از نظر ساختار قیاسی، این نتیجه معتبر است. اما ممکن است فرد به دلیل تجربهٔ شخصی یا وضعیت اقتصادی اطرافش آن را نپذیرد. اینجا نتیجهٔ منطقی معتبر است، اما \eng{prior knowledge} یا تجربهٔ زیسته باعث مقاومت در برابر پذیرش آن می‌شود.

\subsection{مقایسهٔ مدل‌های زبانی و انسان در قیاس}
در بخش پایانی، مقاله‌ای بررسی شد که مدل‌های زبانی را روی فرم‌های صوری \eng{Syllogism} با انسان‌ها مقایسه کرده بود. مدل‌ها با فرم‌هایی مانند \(\text{\eng{All A are B}}\)، \(\text{\eng{Some B are C}}\) و حالت‌های مشابه آزمون شده بودند، نه الزاماً با زبان طبیعی پیچیده.

نتیجهٔ مهم این بود که با بزرگ‌تر شدن مدل‌ها، دقت آن‌ها در برخی موارد از انسان بیشتر می‌شود. اما این به معنای هم‌ترازی کامل با انسان نیست. مدل‌های بزرگ ممکن است پاسخ درست‌تری بدهند، ولی توزیع خطاهایشان شبیه انسان نباشد. یعنی جاهایی که انسان‌ها اشتباه می‌کنند، مدل الزاماً اشتباه نمی‌کند، و جاهایی که مدل خطا دارد، الزاماً با خطاهای انسانی منطبق نیست.

\coursefigure{0.95}{analogy-response-pattern-correlation.png}{همبستگی الگوی پاسخ‌ها در فرم‌های قیاسی مختلف؛ این نمودار برای تحلیل هم‌ترازی خطاها مهم‌تر از دقت خام است، چون شباهت رفتار مدل با الگوهای انسانی را نشان می‌دهد.}

پس عملکرد بهتر از انسان لزوماً به معنای فهم انسانی یا هم‌ترازی با انسان نیست. ممکن است مدل به دلیل حافظهٔ بهتر، دادهٔ آموزشی بیشتر یا دیدن نمونه‌های مشابه در آموزش، پاسخ درست بدهد. این با فهم انسانی یکی نیست.

مدل‌های زبانی معمولاً با متن‌های درست، استدلال‌های رسمی و داده‌های گسترده آموزش می‌بینند. در نتیجه ممکن است فرم‌های درست را زیاد دیده باشند. اما انسان در جهان واقعی با مغالطه، ابهام، تجربهٔ شخصی، بایاس و تصمیم‌گیری سریع روبه‌روست. بنابراین اگر مدل فقط در فرم صوری بهتر عمل کند، هنوز نمی‌توان گفت مثل انسان استدلال می‌کند.

\subsection{حافظه، آمار و استدلال}
یک پرسش پایانی این است که آیا مدل فقط حفظ کرده است یا واقعاً استدلال می‌کند. نکتهٔ مهم این است که «حفظ کردن» تعبیر دقیقی نیست. هم انسان و هم مدل به‌نوعی آماری یاد می‌گیرند. انسان نیز بر اساس تجربه‌های قبلی، \eng{prior}ها و توزیع‌های آماری جهان تصمیم می‌گیرد.

اما آمار مدل با آمار تجربهٔ انسانی یکی نیست. اگر بخواهیم مدل واقعاً به انسان نزدیک‌تر شود، باید توزیع دادهٔ آموزشی و خطاهای آن را به سمت خطاها و الگوهای انسانی تنظیم کنیم.

یکی از راه‌ها \eng{Adversarial Data Augmentation} است: ببینیم مدل در چه نوع نمونه‌هایی شکست می‌خورد، سپس همان نوع نمونه‌ها را بیشتر به داده اضافه کنیم تا مدل روی آن‌ها بهتر شود. در دیتاست \eng{Avicenna} نیز چنین کاری انجام شده بود: فرم‌هایی که مدل در آن‌ها مشکل داشت شناسایی شدند و داده‌های مشابه آن‌ها تقویت شدند. در نتیجه عملکرد مدل بهتر شد، اما این بهبود صرفاً با افزایش تصادفی داده حاصل نشده بود؛ بلکه بر اساس تحلیل خطا انجام شده بود.

جمع‌بندی این بخش این است که برای فهم واقعی مدل، نباید فقط دقت نهایی را سنجید. باید دید آیا مدل در مسیر رسیدن به پاسخ، شبیه انسان عمل می‌کند یا نه. در استدلال قیاسی، انسان‌ها الگوی خطای خاصی دارند و این خطاها تصادفی نیستند؛ بلکه به \eng{prior knowledge}، \eng{System 1}، \eng{System 2}، تجربه، بایاس و ساختار شناختی انسان مربوط‌اند.

اگر مدل در پاسخ نهایی از انسان بهتر باشد اما توزیع خطایش با انسان هم‌تراز نباشد، هنوز نمی‌توان گفت واقعاً با انسان \eng{aligned} است. بنابراین در \eng{Brain Alignment} علاوه بر پاسخ درست، باید به خطاهای انسان، خطاهای مدل، توزیع خطاها، توضیح یا مسیر استدلال انسان، لایه‌های میانی مدل، اثر \eng{prior knowledge} و تفاوت میان \eng{reasoning} صوری و \eng{reasoning} در جهان واقعی توجه کرد.
